% !TeX program = xelatex
\documentclass[aspectratio=169,10pt]{beamer}
\usepackage{amsmath,amssymb,array,booktabs,mathtools}
\usepackage{graphicx,hyperref}
\usepackage{tikz}
% ============================================================
% Beamer Color System — Speculative Decoding black theme
% ============================================================

\providecommand{\beamerthemename}{black}

\RequirePackage{xifthen}

% Keep the public theme selector compatible with the source deck. Its current
% variants intentionally resolve to the same monochrome academic palette.
\newcommand{\usemonochrometheme}{%
  \definecolor{themePrimary}{HTML}{111111}%
  \definecolor{themeSecondary}{HTML}{444444}%
  \definecolor{themeTertiary}{HTML}{000000}%
  \definecolor{themeLight}{HTML}{B8B8B8}%
  \definecolor{themeSilver}{HTML}{E6E6E6}%
  \definecolor{themeBg}{HTML}{FAFAFA}%
  \definecolor{themeBorder}{HTML}{D7D7D7}%
  \definecolor{themeAlert}{HTML}{111111}%
  \definecolor{themeExample}{HTML}{111111}%
}

\ifthenelse{\equal{\beamerthemename}{black}}{\usemonochrometheme}{%
\ifthenelse{\equal{\beamerthemename}{blue}}{\usemonochrometheme}{%
\ifthenelse{\equal{\beamerthemename}{gold}}{\usemonochrometheme}{%
\ifthenelse{\equal{\beamerthemename}{green}}{\usemonochrometheme}{%
\ifthenelse{\equal{\beamerthemename}{purple}}{\usemonochrometheme}{%
  \PackageError{beamer-colors}{Unknown theme '\beamerthemename'.%
    Valid themes: black, blue, gold, green, purple}{}
}}}}}

\setbeamercolor{structure}{fg=themePrimary}
\setbeamercolor{palette primary}{bg=themePrimary,fg=white}
\setbeamercolor{palette secondary}{bg=themeSecondary,fg=white}
\setbeamercolor{palette tertiary}{bg=themeTertiary,fg=white}
\setbeamercolor{palette quaternary}{bg=themePrimary!80!black,fg=white}

\setbeamercolor{title}{fg=themePrimary}
\setbeamercolor{frametitle}{fg=themePrimary}
\setbeamercolor{subtitle}{fg=themeSecondary}
\setbeamercolor{author}{fg=themePrimary!80!black}
\setbeamercolor{institute}{fg=themeSecondary}
\setbeamercolor{date}{fg=themeSecondary}

\setbeamercolor{normal text}{fg=black!84,bg=white}
\setbeamercolor{alerted text}{fg=themeAlert}
\setbeamercolor{example text}{fg=themeExample}

\setbeamercolor{block title}{bg=themePrimary,fg=white}
\setbeamercolor{block body}{bg=themeSilver!40}
\setbeamercolor{block title alerted}{bg=themeAlert,fg=white}
\setbeamercolor{block body alerted}{bg=themeAlert!8}
\setbeamercolor{block title example}{bg=themeExample,fg=white}
\setbeamercolor{block body example}{bg=themeExample!8}

\setbeamercolor{itemize item}{fg=themePrimary}
\setbeamercolor{itemize subitem}{fg=themeSecondary}
\setbeamercolor{itemize subsubitem}{fg=themeSecondary!70}
\setbeamercolor{enumerate item}{fg=themePrimary}
\setbeamercolor{enumerate subitem}{fg=themeSecondary}

\setbeamercolor{section in toc}{fg=themePrimary}
\setbeamercolor{subsection in toc}{fg=themeSecondary}
\setbeamercolor{footnote mark}{fg=themePrimary}
\setbeamercolor{caption name}{fg=themePrimary}

\definecolor{positive}{HTML}{111111}
\definecolor{negative}{gray}{0.42}
\definecolor{emphasis}{HTML}{111111}
\definecolor{neutral}{gray}{0.55}

\newcommand{\pos}[1]{\textcolor{positive}{#1}}
\newcommand{\con}[1]{\textcolor{negative}{#1}}
\newcommand{\HL}[1]{\textcolor{emphasis}{#1}}

\makeatletter
\@ifpackageloaded{hyperref}{%
  \hypersetup{
    colorlinks=true,
    linkcolor=black!75,
    urlcolor=black!75,
    citecolor=black!75,
  }%
}{}
\makeatother

% ============================================================
% Layout: Metropolis — Modern minimal, no external dependencies
% Recreates the metropolis aesthetic using standard beamer.
% ============================================================

\usetheme{default}
\usecolortheme{default}

\newlength{\bsProgressDone}
\newlength{\bsProgressRemain}

% ---- Typography ----
\usefonttheme{professionalfonts}
\setbeamerfont{title}{size=\Large,series=\bfseries}
\setbeamerfont{frametitle}{size=\large,series=\bfseries}
\setbeamerfont{framesubtitle}{size=\normalsize,series=\mdseries}
\setbeamerfont{author}{size=\normalsize}
\setbeamerfont{institute}{size=\small}
\setbeamerfont{date}{size=\small}
\setbeamerfont{section title}{size=\Large,series=\bfseries}
\setbeamerfont{footline}{size=\tiny}

% ---- Remove navigation ----
\setbeamertemplate{navigation symbols}{}

% ---- Frame title: left-aligned, clean ----
\setbeamertemplate{frametitle}{%
  \vskip6pt%
  \usebeamerfont{frametitle}\usebeamercolor[fg]{frametitle}\insertframetitle%
  \ifx\insertframesubtitle\empty\else
    \\[2pt]\usebeamerfont{framesubtitle}\usebeamercolor[fg]{subtitle}\insertframesubtitle%
  \fi
  \vskip4pt%
  {\color{themePrimary}\hrule height 0.45pt}%
  \vskip6pt%
}

% ---- Footline: thin progress bar + page number ----
\setbeamertemplate{footline}{%
  \nointerlineskip%
  \pgfmathsetmacro{\progressfrac}{\insertframenumber/\inserttotalframenumber}%
  \pgfmathsetlength{\bsProgressDone}{\progressfrac\paperwidth}%
  \pgfmathsetlength{\bsProgressRemain}{\paperwidth-\bsProgressDone}%
  \hbox to\paperwidth{%
    {\color{themePrimary!75}\rule{\bsProgressDone}{1pt}}%
    {\color{themeBorder!65}\rule{\bsProgressRemain}{1pt}}%
  }%
  \vskip2pt%
  \hbox to\paperwidth{\hfill\usebeamerfont{footline}\textcolor{black!45}{\insertframenumber\,/\,\inserttotalframenumber}\hspace{8pt}}%
  \vskip4pt%
}

% ---- Title page: minimal ----
\setbeamertemplate{title page}{%
  \vfill
  \begin{flushleft}
    {\usebeamerfont{title}\usebeamercolor[fg]{title}\inserttitle}\\[18pt]
    {\usebeamerfont{author}\textcolor{black!60}{\insertauthor}}\\[4pt]
    {\usebeamerfont{date}\textcolor{black!60}{\insertdate}}
  \end{flushleft}
  \vfill
}

% ---- Section page ----
\AtBeginSection{%
  \begin{frame}[plain,noframenumbering]
    \vfill
    \begin{center}
      {\usebeamerfont{section title}\usebeamercolor[fg]{structure}\insertsectionhead}
    \end{center}
    \vfill
  \end{frame}
}

% ---- Itemize: clean circles ----
\setbeamertemplate{itemize item}{\small\raise1pt\hbox{\textbullet}}
\setbeamertemplate{itemize subitem}{\scriptsize\raise1pt\hbox{--}}
\setbeamertemplate{itemize subsubitem}{\scriptsize\raise1pt\hbox{$\cdot$}}

\author{Skill QA}
\date{5 September 2026}
\newcommand{\slidecite}[1]{%
  \begin{tikzpicture}[remember picture,overlay]
    \node[anchor=south west,align=left,text width=0.92\paperwidth,
      inner sep=0pt,font=\tiny,text=black!38]
      at ([xshift=0.68cm,yshift=0.72cm]current page.south west)
      {\raggedright #1\par};
  \end{tikzpicture}%
}
\newcommand{\norm}[1]{\left\lVert #1\right\rVert}

\title{Why the Midpoint Minimizes Squared Distance}
\hypersetup{pdftitle={Why the Midpoint Minimizes Squared Distance},pdfauthor={Skill QA}}
\begin{document}
\begin{frame}[plain]
  \titlepage
\end{frame}

\begin{frame}{One location determines both squared distances}
  \vspace{0.25cm}
  \[
    \underset{x\in\mathbb{R}^{d}}{\operatorname{minimize}}
    \qquad
    F(x)=\norm{x-a}^{2}+\norm{x-b}^{2}
  \]
  \vspace{0.55cm}
  \begin{itemize}
    \item \(a,b\in\mathbb{R}^{d}\) are fixed. \(d\geq1\).
    \item \(x\) can be any point in the same space.
    \item \(\norm{\cdot}\) is the Euclidean norm. Both terms have equal weight.
  \end{itemize}
  \vspace{0.3cm}
  \begin{center}
    Moving toward one endpoint can move us away from the other.
  \end{center}
  \slidecite{Supplied theorem fixture T1; definitions A1 in source-notes.md.}
\end{frame}

\begin{frame}{The midpoint is the unique minimizer}
  \vspace{0.2cm}
  {\large Theorem}\qquad Define the midpoint \(m=(a+b)/2\).
  \vspace{0.45cm}
  \[
    \underset{x\in\mathbb{R}^{d}}{\operatorname{argmin}}\; F(x)=\{m\},
    \qquad
    F(m)=\frac{\norm{a-b}^{2}}{2}.
  \]
  \vspace{0.5cm}
  \begin{center}
    We will prove the exact gap for every competing \(x\):
    \vspace{0.25cm}
    \[
      F(x)-F(m)=2\norm{x-m}^{2}.
    \]
  \end{center}
  \slidecite{Statement T2 from supplied identity T1; full proof D1-D3 in source-notes.md.}
\end{frame}

\begin{frame}{Centering gives equal and opposite offsets}
  \vspace{0.1cm}
  \[
    y=x-m,\qquad v=\frac{b-a}{2}.
  \]
  \vspace{0.25cm}
  \[
    a=m-v,\qquad b=m+v
  \]
  \vspace{0.35cm}
  \[
    \underbrace{x-a}_{\text{offset from }a}=y+v,
    \qquad
    \underbrace{x-b}_{\text{offset from }b}=y-v.
  \]
  \vspace{0.45cm}
  \begin{center}
    \(y\) measures displacement from the midpoint.\\[0.15cm]
    \(v\) is the fixed half-separation between the endpoints.
  \end{center}
  \slidecite{Algebraic proof D1 in source-notes.md. All equalities hold for any fixed \(a,b\in\mathbb{R}^{d}\).}
\end{frame}

\begin{frame}{The cross terms cancel exactly}
  \vspace{0.1cm}
  \begin{align*}
    \norm{y+v}^{2}
      &=\norm{y}^{2}+2\langle y,v\rangle+\norm{v}^{2},\\[0.2cm]
    \norm{y-v}^{2}
      &=\norm{y}^{2}-2\langle y,v\rangle+\norm{v}^{2}.
  \end{align*}
  \vspace{0.2cm}
  \[
    F(x)=2\norm{y}^{2}+2\norm{v}^{2}
       =2\norm{x-m}^{2}+\frac{\norm{a-b}^{2}}{2}.
  \]
  \vspace{0.35cm}
  \begin{center}
    \(\langle y,v\rangle\) is the Euclidean inner product.\\[0.12cm]
    The endpoint separation contributes only a constant.
  \end{center}
  \slidecite{Full inner-product expansion D2 in source-notes.md, deriving the supplied identity T1.}
\end{frame}

\begin{frame}{A nonnegative gap proves global uniqueness}
  \vspace{0.3cm}
  \[
    F(x)-F(m)=2\norm{x-m}^{2}\geq0
    \qquad\text{for every }x\in\mathbb{R}^{d}.
  \]
  \vspace{0.5cm}
  \[
    F(x)=F(m)
    \quad\Longleftrightarrow\quad
    \norm{x-m}^{2}=0
    \quad\Longleftrightarrow\quad
    x=m.
  \]
  \vspace{0.55cm}
  \begin{center}
    No other point reaches the same value.\\[0.18cm]
    The proof also covers \(a=b\), when \(m=a=b\).
  \end{center}
  \slidecite{Optimization proof D3 in source-notes.md. Positive definiteness supplies the equality condition.}
\end{frame}

\begin{frame}{For 0 and 2, the excess loss is \(2(x-1)^2\)}
  \vspace{0.1cm}
  \[
    a=0,\quad b=2,\quad m=1,\qquad
    F(x)=x^2+(x-2)^2=2(x-1)^2+2.
  \]
  \vspace{0.5cm}
  \begin{center}
    \renewcommand{\arraystretch}{1.5}
    \begin{tabular}{lccccc}
      \toprule
      \(x\) & 0 & 0.5 & \(\mathbf{1}\) & 1.5 & 2\\
      \midrule
      \(F(x)\) & 4 & 2.5 & \(\mathbf{2}\) & 2.5 & 4\\
      \bottomrule
    \end{tabular}
  \end{center}
  \vspace{0.5cm}
  \begin{center}
    Equal displacements from 1 incur equal excess loss.
  \end{center}
  \slidecite{Worked example E1 in source-notes.md. This example illustrates the proved theorem.}
\end{frame}

\begin{frame}{The guarantee depends on squared Euclidean geometry}
  \vspace{0.2cm}
  Without the square, uniqueness already fails on the real line:
  \vspace{0.2cm}
  \[
    |x|+|x-2|=2
    \qquad\text{for every }x\in[0,2].
  \]
  \vspace{0.35cm}
  \begin{itemize}
    \item An arbitrary metric need not define the vector midpoint or satisfy the identity.
    \item The theorem assumes equal weights and an unconstrained Euclidean domain.
  \end{itemize}
  \vspace{0.35cm}
  \begin{center}
    In the stated setting, the exact squared gap proves one global minimizer.
  \end{center}
  \slidecite{Scope contrast S1 in source-notes.md. The unsquared example follows from the triangle inequality.}
\end{frame}
\end{document}
